\documentclass[11pt]{article}
\usepackage{amsmath,amssymb,amsthm,mathtools,hyperref}

\begin{document}
\section*{Response to novel environments}
\label{prob:182}
\label{rank:94}
\par\smallskip

\noindent\textbf{Domain.} Organismal ecology; novel environments. \quad \textbf{Retained status/source context.} Open constrained extrapolation problem. Prediction outside the observed environment is not statistically identifiable without additional mechanistic assumptions.
\paragraph{Definitions and assumptions.} For a plant or ectotherm genotype $g$, let physiological state $x_t$ contain measured energy reserves, damage and acclimation state. Environment $e_t$ records temperature, moisture and resource concentration in stated units. Candidate dynamics are $x_{t+1}=f_\theta(x_t,e_t)$ and reproduction $W_t=h_\theta(x_t,e_t)$, with $\theta$ in a bounded mechanistic parameter set calibrated by controlled stresses. New environments belong to a declared scenario set $\mathcal E_{\mathrm{new}}$, which may extend beyond training support but respects physical bounds. The same physiological state variables must be meaningful throughout the scenario range. If a novel stress introduces a previously absent damage mechanism, the model class requires expansion and previously narrow bounds may no longer have their stated coverage.
\paragraph{Problem statement.} Construct scenario-specific fitness and survival bounds,
\[
 W_-(e_{0:T})=\inf_{\theta\in\mathcal C(D)}\sum_{t=0}^T W_t(\theta),
 \qquad W_+(e_{0:T})=\sup_{\theta\in\mathcal C(D)}\sum_{t=0}^T W_t(\theta),
\]
and identify the additional experiments needed to reduce $W_+-W_-$ below tolerance $\varepsilon$ for specified novel stress sequences. Here $\mathcal C(D)$ is a calibrated parameter confidence set.
\paragraph{Additional binding assumptions and resolution criterion.} Specify survival dynamics and whether the reproductive total counts only living individuals or descendants, with a fixed census convention. A confidence set for parameters gives scenario bounds only conditional on the model remaining correct in the novel environment. Separate extrapolation-model discrepancy from parameter uncertainty, and declare the admissible experimental designs and cost before asking which experiments achieve the target interval width. The displayed envelopes bound the reproductive total; survival bounds remain a jointly required output with a declared survival endpoint and readout. Declare the units and census timing of $W_t$, any conversion from a rate to interval reproduction, and whether the target is deterministic, an expected output or a realized random outcome; descendant accounting requires dynamics sufficient to define that total. For $\mathcal C(D)$, declare the calibration and observation law, confidence level and scope of parameter coverage, and how initial uncertainty and biological process variation enter the target. Parameter coverage alone does not provide predictive coverage for realized outcomes or coverage of extrapolation-model discrepancy. For additional experiments, declare the available observations and feedback, budget, confidence-set update and any stopping rule, together with the quantifiers over parameters, future data and specified scenarios in the width requirement. Whether that requirement is worst-case, probabilistic or in expectation, and whether it is pointwise or simultaneous across scenarios, remains to be specified; no one convention is selected here. The recursions and readouts must be well-defined on the admitted trajectories. Bounded parameters alone do not ensure finite envelopes; an infinite interval width does not meet the stated positive finite tolerance. No attainment of the infimum or supremum is assumed.

\paragraph{Alternative formulations.} (i) Freeze the genotype and predict physiological acclimation/plasticity over one lifetime. (ii) Allow heritable adaptation across generations under bounded mutation and migration, predicting evolutionary rescue probability instead of individual fitness. These are different processes and timescales.
\paragraph{Obstructions and partial progress.} The volume names response to complex novel environments explicitly. Genomic forecast-design work also emphasizes experimental validation. Multiple physiological models can agree within training conditions yet disagree under combined extremes, so generic machine-learning accuracy cannot certify extrapolation. The research target is experimentally constrained mechanistic prediction with honest scenario bounds, including path dependence and interactions among stressors, rather than unrestricted prediction for every previously unobserved environment.
\paragraph{Sources (supplied provenance; not all rechecked).}
\begin{itemize}
\item Dobson, Holt and Tilman (eds.), Unsolved Problems in Ecology (2020). \url{https://www.jstor.org/stable/j.ctvs9fh2n}. Supplied book; contents and chapter previews checked, full text unavailable. Historical agenda, not present-day certification. Relevant listed chapter: Snell-Rood and Kobiela, How Will Organisms Respond to Complex, Novel Environments?.
\item Lotterhos, Principles in experimental design for evaluating genomic forecasts (2024). \url{https://besjournals.onlinelibrary.wiley.com/doi/full/10.1111/2041-210X.14379}. Methodological primary source on validation and reciprocal transplants.
\end{itemize}


\section*{Binding catalogue conventions}
This edition revises both supplied catalogues and preserves original problem numbers 1--330. It adds one globally unique label \verb|\label{rank:x}| per problem, where smaller $x$ means greater provisional difficulty. A rank label is an anchor at that problem's section: \verb|\ref{rank:x}| returns the section number, not the literal difficulty number. The visible \texttt{rank:x} text and the index give the difficulty number itself.

The entries are \emph{parameterized mathematical research schemes}, including empirical prediction, causal identification, historical reconstruction, theory and engineering. They are not 330 certified unsolved mathematical conjectures. A rigorous instance fixes the objects and restrictions below; changing those objects changes the question. Established identities, finite calculations, non-identification constructions and solved special cases are not being posed anew as open. Some residual tasks concern informative bounds, biological realization or external validation rather than a new theorem.

Every main definition and problem statement received a mathematical editorial review. Each entry now has an additional binding-conditions paragraph; direct corrections are separately recorded in the accompanying audit. The alternative formulations inherit the applicable conventions, but are different targets unless equivalence is proved. Where an old broad phrase conflicts with an explicit restriction below it, the explicit restriction controls.

\textbf{Literature boundary.} Thirteen selective primary-literature checks cover eighteen entries in this revision; their access levels and scope are recorded explicitly. The supplied bibliography and its earlier source-access descriptions are retained as provenance, not represented as newly verified. Neither this review nor a source's description proves the absence of later solutions. The date of this edition is 13 October 2026; there is no claim of exhaustive current-status certification for all 330 entries.

\subsection*{Binding conventions for a concrete instance}
\paragraph{Objects, quantifiers and non-vacuity.}
An instance declares measurable state, input, output and parameter spaces; a fixed candidate family; admissible initial laws, interventions and environments; a finite horizon or specified asymptotic regime; observation and sampling laws; and performance metrics, positive normalization scales and positive tolerances. These are chosen before fitting or testing, not adjusted after seeing failures. Shared parameters must genuinely be shared across the declared validation domain. Finite dimension alone does not restrict arbitrary real-number encodings: specify a finite library, bounded description length, or regularity/complexity bounds. Exclude constant or no-signal solutions through a fixed informative-output or baseline-improvement requirement where appropriate. A universally quantified result cannot be established by choosing one trivial favorable instance.

\paragraph{Probability laws and data.}
Every expectation or probability refers to the specified stochastic law, including initial uncertainty, process noise, observation noise and sampling as appropriate. A deterministic equation is not a probability model without a declared random input or initial law. ``Observed'' targets in population-level formulas mean population observables $g_j^*$, not noisy estimates asserted to be exact. Let $\widetilde g_j$ be independent validation estimates and suppose
\[
\mathbb{P}\{\lvert\widetilde g_j-g_j^*\rvert\le r_j\text{ for every }j\}\ge1-\alpha.
\]
Then an uncertainty-aware sufficient certification rule for positive target tolerances $\varepsilon_j$ is
\[
\max_j\frac{|\widehat g_j-\widetilde g_j|+r_j}{\varepsilon_j}\le1.
\]
On that simultaneous confidence event it implies every population error is at most its tolerance, by the triangle inequality. Vector/norm versions use corresponding valid simultaneous bounds. Training uncertainty must also be included when the claim averages over repeated training data. Evaluation on finitely many conditions does not establish a continuous-domain supremum without covering and regularity bounds.

\paragraph{Dynamics and well-posedness.}
For ODEs and SDEs give coefficient domains, initial laws and existence/uniqueness conditions in the intended solution sense, or keep the entire admitted solution set. For stochastic jump models specify nonnegative off-diagonal rates, diagonal normalization, rate orientation and non-explosion. For PDEs specify spatial domains, boundary/initial conditions, conservation or source terms, and a weak or classical solution class. Moving-domain concentration equations include dilution/transport terms. A stationary law is used only for a time-homogeneous augmented model under existence assumptions; a stationary approximation to a driven system must be justified. Absorbing extinction states are retained. Survival-conditioned, quasi-stationary and unconditional laws are distinct and must not be substituted for one another.

\paragraph{Identification and causal meaning.}
Two parameter values are observationally equivalent when they induce the same full distribution under all allowed observation/intervention designs, not merely the same likelihood value for one sample. A target is identified only if constant on the equivalence class. Gauge symmetries, label switching and reference conventions are quotiented or fixed. Finite-data compatible sets use a declared confidence, likelihood or discrepancy rule; an empty set signals model rejection. Bayesian credible sets are not automatically confidence sets. A prior or regularizer does not create data identification. Bounds intended to resolve a question must be informative at a predeclared width or decision threshold.

A causal contrast uses well-defined intervention versions and a stated causal model. Random assignment, or explicit consistency, overlap and exchangeability assumptions, is required for causal interpretation from observations. Interference and time-varying confounding are represented when present. Fixing a mediator by intervention differs from conditioning on its observed value. Marginal potential-outcome laws do not determine individual cross-world effects or an arbitrary joint coupling. Ethically admissible designs, rather than unconstrained interventions, define the experimental menu.

\paragraph{Metrics, units, optimization and limits.}
Specify the sampling or reference measure for densities and integrals, the log base, and dimensionless arguments of logarithms. Total variation is $\sup_A|P(A)-Q(A)|$; Wasserstein distances require the declared ground metric and appropriate moment conditions; KL divergence uses its stated absolute-continuity convention. All normed readouts are scaled consistently, and error probabilities use $0<\alpha<1$. A ratio has a positive denominator or an explicit undefined/zero convention. Thermodynamic entropy production is nonnegative in expectation under the relevant hypotheses; individual stochastic realizations need not be. An infimum is not an attained minimum without compactness/coercivity and continuity or another existence argument. Fix the order of large-population, long-time and vanishing-mutation limits. Stability claims declare perturbation space and symmetries.

\subsection*{How the provisional difficulty order was assigned}
Four editorial components are scored from 1 (lower burden) to 5 (higher burden): (a) formal or inferential difficulty, (b) access to evidence and identification burden, (c) mechanistic/computational integration, and (d) independent validation and transfer. Their unweighted total orders the entries. For selected entries in the highest bands, an explicit editorial priority breaks equal totals; all other ties are broken by original problem number. Equal totals are \emph{not} evidence of reliably different true difficulties. Scores are judgments, not measurements, completion-time estimates or consensus rankings.

The object ranked is an informative, nondegenerate resolution of the \emph{scoped} scientific target, with the stated budget and validation domain. A sharp useful impossibility result can qualify; an already known identity, vacuous bound, arbitrary favorable model or lookup table cannot. Rank 330 is merely the least difficult relative assessment in this catalogue, not a claim of an easy biological solution. Broad historical titles do not automatically outrank precisely formulated mechanistic tasks.



\end{document}
